\section{Comparison with Greedy Merging}
\label{ssec:greedy-comparison}

We compare our method with greedy merging under the same objective.
To construct parent domains from consecutive atomic TADs, we initialize the ordered list of leaf ranges as $\mathcal P_0=([0,1),\ldots,[m-1,m))$.
At iteration $t\in\{0,\ldots,m-2\}$, let $s=m-t$ be the number of current domains and write $\mathcal P_t=(I_0,\ldots,I_{s-1})$, where $I_r=[a_r,a_{r+1})$ for all $r\in\{0,\ldots,s-1\}$, with integer leaf indices $0=a_0<\cdots<a_s=m$.
Merging $I_r$ and $I_{r+1}$ adds $W_{a_r,a_{r+2}}$ to the sum of parent-domain costs.
We choose the merge with the smallest added cost:
\[
 r^\star=\min\operatorname*{arg\,min}_{r\in\{0,\ldots,s-2\}}W_{a_r,a_{r+2}}.
\]
The outer minimum selects the leftmost pair when several pairs have the same cost.
We add a parent node representing $[a_{r^\star},a_{r^\star+2})$, with $I_{r^\star}$ and $I_{r^\star+1}$ as its two children.
We update the list by
\[
 \mathcal P_{t+1}=(I_0,\ldots,I_{r^\star-1},\ I_{r^\star}\cup I_{r^\star+1},\ I_{r^\star+2},\ldots,I_{s-1}),
\]
where the prefix or suffix is omitted when the selected pair is at an end of the list.
After $m-1$ merges, $\mathcal P_{m-1}=([0,m))$, and the recorded parent--child relationships define a full ordered binary tree.
We fix the leaves, interval costs, fractions of parent domains retained, length caps, and ranking rules for both methods.
Each of the seven combinations of cell line and resolution includes all 22 autosomes.

In Table~\ref{tab:greedy-point}, greedy merging has higher adjusted $R^2$ in all seven combinations.
Our method has higher intra-inter difference in six combinations and a larger fraction of parent domains with CTCF peaks near both endpoints in one.
Exact minimization of the proposed objective therefore does not imply better scores on every evaluation metric.

\begin{table}[htbp]
\centering
\caption{Values are Greedy / Ours, with the best value in each comparison shown in bold. Intra-inter difference values are multiplied by $10^4$ for 5 kb inputs only. CTCF gives the percentage of parent domains with peaks within $\pm10$ kb of both endpoints.}
\label{tab:greedy-point}
\footnotesize
\setlength{\tabcolsep}{4pt}
\begin{tabular}{llrrr}
\toprule
Cell & Res. & Adjusted $R^2$ & Intra-inter difference & CTCF\\
\midrule
GM12878 & 25 kb & \bestvalue{0.1837} / 0.1822 & 44.244 / \bestvalue{44.797} & \bestvalue{30.47} / 30.02\\
GM12878 & 10 kb & \bestvalue{0.0827} / 0.0819 & 15.643 / \bestvalue{16.108} & \bestvalue{42.32} / 41.30\\
K562 & 25 kb & \bestvalue{0.0919} / 0.0913 & 6.726 / \bestvalue{6.797} & \bestvalue{39.74} / 38.08\\
IMR90 & 25 kb & \bestvalue{0.2323} / 0.2298 & 40.568 / \bestvalue{40.943} & 35.33 / \bestvalue{37.24}\\
GM12878 & 5 kb & \bestvalue{0.1411} / 0.1392 & \bestvalue{57.108} / 56.962 & \bestvalue{50.50} / 49.09\\
K562 & 5 kb & \bestvalue{0.0717} / 0.0703 & 65.831 / \bestvalue{66.779} & \bestvalue{60.28} / 58.37\\
IMR90 & 5 kb & \bestvalue{0.1108} / 0.1105 & 87.281 / \bestvalue{88.431} & \bestvalue{51.05} / 50.69\\
\bottomrule
\end{tabular}
\end{table}
